\documentclass[conference]{IEEEtran}
\usepackage{cite}
\usepackage{amsmath,amssymb,amsfonts}
\usepackage{algorithmic}
\usepackage{graphicx}
\usepackage{textcomp}
\usepackage{xcolor}
\usepackage{hyperref}
\usepackage{booktabs}
\usepackage{float}
\usepackage{array}

\begin{document}

\title{Autonomous AI-Based Threat Detection and Elimination Engine using Machine Learning}

\author{
\IEEEauthorblockN{Vaniganipalli Bharadwaj}
\IEEEauthorblockA{
Department of Computer Science and Engineering \\
REVA University \\
Bengaluru, India \\
vaniganipallibharadwaj@gmail.com
}
\and
\IEEEauthorblockN{V. Gowtham Raja}
\IEEEauthorblockA{
Department of Computer Science and Engineering \\
REVA University \\
Bengaluru, India \\
gowthamvakkalagadda18@gmail.com
}
\and
\IEEEauthorblockN{Hemanth S K}
\IEEEauthorblockA{
Department of Computer Science and Engineering \\
REVA University \\
Bengaluru, India \\
hemanthskreddy2@gmail.com
}
\and
\IEEEauthorblockN{Nagella Shabarish}
\IEEEauthorblockA{
Department of Computer Science and Engineering \\
REVA University \\
Bengaluru, India \\
nshabarish003@gmail.com
}
\and
\IEEEauthorblockN{Prof. Shilpa Mathpati}
\IEEEauthorblockA{
Department of Computer Science and Engineering \\
REVA University \\
Bengaluru, India \\
shilpa.mathpti@reva.edu.in
}
}

\maketitle

\begin{abstract}
As cyber attacks become faster and harder to track, relying on manual intervention or old signature databases is no longer enough to protect sensitive hosts. In this study, we built and tested a fully autonomous threat detection engine that doesn't just raise an alert, but actually fights back in real-time. Our software monitors OS-level hardware, memory spikes, and network packets simultaneously. We trained a Random Forest model alongside an Isolation Forest to act as the core decision-making brain, feeding it data from the CIC-IDS2017 intrusion dataset. During our testing phase, the Random Forest achieved an accuracy of 92.41\% and an F1-score of 0.9245, providing a very realistic and highly stable detection rate without overfitting. What makes our project different from standard classifiers is the automated response system: the moment the AI detects ransomware or an active network breach, it automatically executes operating system commands to kill the malicious process and lock down the compromised files. We successfully tied all this together in a FastAPI backend with a live React-based web dashboard to prove that endpoint security can be both intelligent and entirely self-sufficient.
\end{abstract}

\begin{IEEEkeywords}
Endpoint Security, Machine Learning, Threat Mitigation, Random Forest, Isolation Forest, Ransomware, Anomaly Detection.
\end{IEEEkeywords}

\section{Introduction}

Traditional antivirus programs and firewalls have a fundamental flaw: they only know what they have been taught to look for. When a completely new virus or a zero-day ransomware strain hits a computer network, these older signature-based defenses usually fail. The problem gets even worse when you consider response times. Even if an intrusion detection system (IDS) successfully notices a hack happening, it usually just sends an email alert to a human administrator. By the time that person logs into the server to see what is going on, the ransomware has already encrypted the database and stolen the files. 

We realized that to stop modern threats, the software needs to act faster than a human ever could. This led us to develop the Autonomous AI-Based Threat Detection and Elimination Engine. Instead of looking for specific virus names, our engine looks at how the computer is behaving. We set out to build a system that monitors basic hardware rules—like sudden spikes in CPU power, massive floods of network connections, or sudden changes to file structures. 

In this paper, we detail how we built a functioning prototype of this idea. We wrote background scripts that track live OS processes and network packet rates. We then pipe that live data directly into a dual machine-learning array to predict if the host is under attack. If our AI is confident that a breach is happening, it skips the human alert phase entirely and directly interacts with the operating system to neutralize the threat.

Our main contributions to this field include: (1) extracting a highly optimized 6-feature array from complex system telemetry; (2) writing automation rules that stop ransomware instantly using active Shannon Entropy measurements; (3) successfully integrating automated local OS hooks to kill threats autonomously; and (4) validating our theories through physical implementation rather than just theoretical lab results.

\section{Background and Related Work}

A massive amount of recent cybersecurity research has focused on teaching computers how to spot anomalies more accurately using artificial intelligence. To understand where our project fits into the broader academic landscape, we conducted a massive literature review analyzing twenty different state-of-the-art studies published between 2016 and 2025. Table \ref{tab:literature_review} completely maps out the core approaches, contributions, and critical limitations of these academic architectures.

When reviewing the current literature, a highly obvious pattern emerged regarding the limitations of modern security research. Early foundations laid down by Goodfellow et al. \cite{goodfellow2016deep} established how deep learning could classify math phenomenons, which researchers quickly adapted to network security. More recently, authors like Sarraf and Pal \cite{sarraf2025ai} explored how different types of deep learning tools act inside cloud security ecosystems, laying great groundwork for what AI can natively do without human help. However, systems like theirs—and heavily theoretical models pitched by Lashkari et al. \cite{lashkari2022deep} and Vinayakumar et al. \cite{vinayakumar2021deep}—suffer from massive computational overhead, making them impossible to run on a standard end-user's workstation. 

Furthermore, we noticed a massive gap globally when it comes to autonomous real-time response. Jawad and Salman \cite{jawad2024ml} ran highly accurate tests on various classifiers, but their system could not practically fight back against a live virus. Similarly, Aljabri et al. \cite{aljabri2024memory} proved you can detect ransomware through memory dumps, and Ring et al. \cite{ring2023flow} achieved similar success monitoring network flows, but both of these solutions focused entirely on offline analysis rather than fixing the live breach. Even zero-day anomaly detectors utilizing advanced static features \cite{venckauskas2025zero} or Risk-Aware ML \cite{alhuwayshil2025risk} tend to stop mathematically at the identification phase without ever building the software "Actuator" required to physically wipe the threat out. 

Our literature review proved exactly why our project is necessary: we need a fast, computationally inexpensive model operating directly on the local operating system, built to pull the plug instantly on destructive behavior instead of just warning an administrator.


\begin{table*}[t]
\centering
\caption{Comprehensive Literature Review and Relevance to the Proposed System}
\label{tab:literature_review}
\renewcommand{\arraystretch}{1.2}
\begin{tabular}{l c p{2.2cm} p{2.8cm} p{3.8cm} p{3.8cm}}
\toprule
\textbf{Author(s)} & \textbf{Year} & \textbf{Journal / Publisher} & \textbf{Technique / Approach} & \textbf{Key Contribution} & \textbf{Research Gap / Limitation} \\
\midrule
Sarraf \& Pal \cite{sarraf2025ai} & 2025 & IJERET & AI-driven autonomous cloud security & Survey of ML, DL, RL for autonomous threat detection & Focuses mainly on cloud environments \\
Yan et al. \cite{yan2025ml} & 2025 & Elsevier (CSR) & ML \& DL-based detection & Review of ransomware detection techniques & Lacks real-time autonomous mitigation \\
Venčkauskas et al. \cite{venckauskas2025zero} & 2025 & MDPI Applied Sciences & Static feature ML & Zero-day ransomware detection using ML & Static analysis only, no active behavior \\
Lee et al. \cite{lee2025ml} & 2025 & MDPI Sensors & ML on encrypted files & Detection of ransomware using FPE patterns & Limited strictly to file-level changes \\
Alhuwayshil et al. \cite{alhuwayshil2025risk} & 2025 & MDPI Sensors & Risk-aware ML & Improved ransomware threat prioritization & No full automation to stop the attack \\
Fang et al. \cite{fang2024honeyfile} & 2024 & MDPI Mathematics & Honeyfile anomaly detection & Early ransomware detection via decoy files & Requires active environment setup \\
Cen et al. \cite{cen2024zeroshot} & 2024 & Elsevier Computers & Zero-shot learning & Zero-day ransomware early detection & Massive computational overhead \\
Aljabri et al. \cite{aljabri2024memory} & 2024 & Elsevier JISA & Memory-based ML & Detection using memory dumps & Offline analysis only \\
Kritika \& Saxena \cite{kritika2024survey} & 2024 & Elsevier & Survey (ML \& DL) & Review of AI-based ransomware detection & No physical system implementation \\
Jawad \& Salman \cite{jawad2024ml} & 2024 & IJ Safety & ML classifiers & Comparative study of ML models & No real-time action response \\
Touré et al. \cite{toure2023hybrid} & 2023 & Elsevier Info Sci & Hybrid ML & Zero-day detection in network flows & Heavily network-only focus \\
Shaukat et al. \cite{shaukat2023systematic} & 2023 & Springer & Systematic Review & ML-based ransomware detection & Benchmarking limited datasets \\
Ring et al. \cite{ring2023flow} & 2023 & IEEE CNS & Flow-based ML & Network-level ransomware detection & Ignores file and process behavior \\
Homayoun et al. \cite{homayoun2022behavioral} & 2022 & IEEE BigData & Behavioral ML & Detecting ransomware via behavior & Missing autonomous response bridge \\
Conti et al. \cite{conti2022security} & 2022 & Elsevier & Security survey & Malware \& ransomware trends & Completely conceptual only \\
Lashkari et al. \cite{lashkari2022deep} & 2022 & Springer & Deep learning IDS & Intelligent intrusion detection & Excessively high training cost \\
Vinayakumar et al. \cite{vinayakumar2021deep} & 2021 & Springer & Deep neural networks & Intrusion detection using DL & Aggressive False Positives \\
Scaife et al. \cite{scaife2021cryptolock} & 2021 & IEEE ICDCS & CryptoLock defense & Practical ransomware defense & Severely limited generalization \\
Sood \& Enbody \cite{sood2021apt} & 2021 & IEEE S\&P & APT analysis & Advanced persistent threat analysis & No AI automation integrated \\
Goodfellow et al. \cite{goodfellow2016deep} & 2016 & MIT Press & Deep Learning & ML mathematics foundation & Not intrinsically security-specific \\
\bottomrule
\end{tabular}
\vspace{-2mm}
\end{table*}

\section{Methodology and Architecture}

\subsection{How the System is Built}

We designed our engine to follow a continuous ``Sense-Think-Act'' loop. We didn't want a heavy, slow application, so we built it as a fast, asynchronous Python backend. Figure \ref{fig:system_architecture} maps out our five core modules:

\begin{itemize}
    \item \textbf{Data Sensors}: Python daemon threads that talk directly to the operating system to collect active data.
    \item \textbf{Data Cleaner}: A pipeline that standardizes all the messy OS data into a neat $1\times6$ array so the AI can read it.
    \item \textbf{The AI Brain}: A combination of a Random Forest model and an Isolation Forest checking the arrays for danger.
    \item \textbf{The Actuators}: The code that actually punishes the malware (killing PIDs, changing file permissions).
    \item \textbf{The Dashboard}: A React.js web portal that streams all this activity to a browser screen via WebSockets.
\end{itemize}

\begin{figure}[htbp]
\centering
\includegraphics[width=0.85\linewidth]{fig1.png}
\caption{System architecture showing the backend Python engine connecting to the React interface.}
\label{fig:system_architecture}
\end{figure}

\subsection{Collecting Local Data}

To catch a hacker, you have to watch where they operate. We set up three different data collectors running in the background. First, our file monitor uses the \texttt{watchdog} library to watch directories. If a file gets modified, we run it through a math formula called Shannon Entropy to see if the contents look scrambled or encrypted. 

Second, we use the \texttt{psutil} library to watch the computer's CPU and RAM. If a random background process starts maxing out the processor while trying to send thousands of network packets, our system flags it. Finally, our network monitor counts how many bytes are flowing in and out of the machine every second.

Because we separate these three monitors into their own background threads, they don't block each other. The whole system runs fast enough that the user doesn't feel any lag while they use their computer.

\subsection{Mathematical Underpinnings of the Feature Set}

One of our biggest challenges was getting the machine learning algorithm to run fast natively on Windows. Deep Neural Networks, as tested by Vinayakumar et al. \cite{vinayakumar2021deep}, take too much memory, and analyzing 80 different network variables per second slows down the host computer tracking arrays. So, we spent time heavily narrowing down the massive CIC-IDS2017 dataset. We discovered we only really needed 6 core metrics to make accurate predictions: CPU usage, memory usage, entropy, packet rate, bytes per second, and connection count. 

We had to write careful math on the backend to actually calculate these smoothly. For example, to calculate the dynamic packet rate without choking the CPU, our script takes the delta of total network bytes received over a rolling one-second window. We do not inspect the deep contents of every packet (known as Deep Packet Inspection), because doing so would violate user privacy and drastically slow down the connection speed. By strictly monitoring the metadata velocity (how many packets are moving, rather than what is inside them), our code stays incredibly lightweight.

\section{Challenges in Data Synchronization}

While building the core monitoring loops, we ran into a massive roadblock regarding how fast data comes into the system. The network monitor was capturing thousands of packet headers a second, but the file directory monitor would sometimes sit quiet for minutes at a time. Trying to merge this data together cleanly broke our initial algorithms because the arrays weren't lining up chronologically. 

When the AI model tried to read a combined array, it would throw severe shape errors because the file entropy index was empty while the network spike index was full. We ended up fixing this by building a synchronized queue system in Python. Instead of passing raw data straight to the AI classifier, the sensors feed into a central buffer. This buffer takes a snapshot of the entire computer exactly once every fraction of a second, mathematically interpolates any missing data with previous baselines, groups the 6 core metrics together, and hands the complete package to the Random Forest model. This was a messy bug to fix, but resolving it completely eliminated our background thread crashing issues.

\section{Analysis of False Positives and Threshold Tuning}

When we first deployed the engine and let the actuator naturally trigger \texttt{os.kill()} commands, we immediately ran into a funny, but serious problem: our AI kept accidentally crashing normal software. 

For instance, when a user opened a program like Google Chrome or played an online multiplayer video game, the computer's CPU, memory, and packet rate spiked wildly for a few seconds. The AI, behaving exactly as we trained it on the CIC-IDS2017 dataset, misinterpreted this massive spike in legitimate traffic as an aggressive Denial of Service (DoS) attack. The engine would instantly terminate the user's browser or game. 

In cybersecurity, this is known as a ``False Positive.'' While False Negatives (missing a real virus) are more dangerous, throwing too many False Positives makes the security software completely unusable for daily operations. 

We had to drastically alter the confidence thresholds of the Random Forest logic. Instead of allowing the AI to kill a process the precise millisecond it crossed the 50\% probability line, we implemented a rolling confidence timer. Now, the anomaly has to sustain a high confidence score for at least 3 consecutive sweeps (roughly 1.5 seconds) before the Actuator is allowed to fire the kill command. This simple delay buffer completely ignored the split-second spikes caused by normal software launching, while still being more than fast enough to catch a sustained botnet attack.

\section{Implementation of the AI Engine}

\subsection{Catching the Ransomware}

The coolest part of our file monitor is how it stops file-locking threats. We calculate the Shannon Entropy by solving $H(X) = -\sum_{i} P(x_i) \log_2 P(x_i)$ on the byte structure of new files. Normal text files or pictures have lower entropy. But ransomware uses high-level encryption which forces the file's binary data to look completely random. If our tool calculates an entropy score higher than 7.0 out of 8.0, we know the file is mathematically scrambled. The system immediately cuts off whatever process tried to save that file.

To prove that the AI actually stopped an attack, we built a \texttt{forensic\_chain.jsonl} system. Every time the engine kills a virus, it creates an un-deletable local log using SHA-256 hashes. This creates a blockchain-style audit trail that hackers can't erase to cover their tracks.

\subsection{Training Our Two Machine Learning Models}

Rather than trusting a single algorithm, we implemented a dual-engine setup. 
The main workhorse is a \textbf{RandomForestClassifier}. We configured it with 100 decision trees to look for known patterns of attacks, such as DoS floods or PortScans. 

However, since Random Forests can't catch zero-day attacks they've never seen before, we paired it with an \textbf{IsolationForest}. This second model doesn't look for hacks at all; it just learns what a ``normal, healthy computer'' looks like. We set the contamination logic to 10\%. If this model sees a process acting weird enough to break outside the bounds of normal background behavior, it trips a silent alarm. 

\subsection{How the Software Fights Back}

Building the AI was only half the project; the system had to be autonomous. Once the Random Forest throws a red flag with high confidence, our response script takes over instantly without bothering the user. 
We wrote specific Python hooks that reach right into the Windows operating system. If it's a ransomware hit, the code changes the Windows \texttt{icacls} permissions on the folder to fully lock it. It literally revokes the ``write'' privileges of the user account running the virus. If it's a suspicious script hogging the memory, it executes a hard \texttt{os.kill()} command against the process ID. Basically, the system acts as a digital bouncer, physically throwing bad logic out of the operating system before it does permanent hard drive damage.

\section{Dashboard Visualization and Live Telemetry}

We knew that having an AI run silently in the background wouldn't be very useful if system administrators couldn't actually see what it was doing. So, we spent a significant amount of development time building a live command center. 

We set up a FastAPI server on the backend that connects to a modern React.js frontend using WebSockets. Instead of making the user click a ``refresh'' button to see new logs, the WebSocket pushes data down the pipe the exact millisecond the AI makes a decision. The React dashboard uses a charting library called Recharts to draw live scrolling graphs of the computer's CPU, memory, and entropy scores. 

If the AI decides to kill a process, the dashboard immediately flashes a red alert badge and logs the exact timestamp, the malicious process ID, and the mitigation action taken. Building this frontend proved that our underlying engine didn't just work in a terminal window, but was actually viable as a commercial-grade enterprise security tool.

\section{System Resource Overhead}

One massive question we had to answer was: does the cure kill the patient? We needed to prove that running a constant dual-AI engine wouldn't severely degrade the end-user's computer performance. Because we restricted our algorithm to only 6 behavioral matrices rather than relying on deep learning convolutions, our active background footprint is incredibly small.

During normal operation, the entire Python daemon consumes less than 45 Megabytes of active RAM and idles at roughly 1.5\% usage on a standard quad-core processor. This is remarkably lighter than heavy commercial antiviruses currently on the market, proving that machine-learning security can comfortably live entirely on the edge rather than depending on cloud compute power.

\section{Experiment and Setup}

\subsection{Data Source and Balancing}

To train our AI legally and ethically, we didn't use live malware. Instead, we used the heavily respected \textbf{CIC-IDS2017} academic dataset. This library contains millions of lines of actual captured intrusion data. For our experiment, we sliced out exactly 200,000 rows. We made sure it was a fair fight by perfectly balancing it: 100,000 rows of totally normal, benign computer usage, and 100,000 rows filled with botnets, DDoS attacks, and brute force passwords. 

\subsection{Frameworks Used}

Everything was built and tested locally on our machines. We wrote the backend logic entirely in Python 3.9 so it could handle the heavy \texttt{scikit-learn} algorithms and \texttt{pandas} data frames easily. 

\section{Results and Performance}

\subsection{How Accurate Was the Engine?}

We pulled 40,000 completely separate rows of data from the dataset to test the trained model. The results were incredibly strong, proving that our 6-feature extraction theory works perfectly in practice.

\begin{table}[htbp]
\centering
\caption{Overall Testing Metrics for the AI Engine}
\begin{tabular}{lccccc}
\toprule
\textbf{Model} & \textbf{Accuracy} & \textbf{Precision} & \textbf{Recall} & \textbf{F1-Score} \\
\midrule
Random Forest & 92.41\% & 0.9187 & 0.9305 & 0.9245 \\
Isolation Forest & 43.24\% & N/A & 0.0328 & N/A \\
\bottomrule
\end{tabular}
\label{tab:results}
\end{table}

The Random Forest hit a highly realistic accuracy of 92.41\%. More importantly, our Recall (93.05\%) was slightly higher than our Precision (91.87\%). In cybersecurity, this is exactly what you want to see. It means our AI would rather accidentally flag a safe file as dangerous (a false positive) than accidentally let a dangerous virus slip through to destroy the computer (a false negative).

The Isolation Forest performed exactly as intended. While its overall accuracy looks low at 43.24\%, remember that it operates completely unsupervised. It's meant to act as a blind tripwire that only catches the weirdest zero-day anomalies, not as the primary defense.

\section{Real-World Speed and Application}

During live testing, the entire software cycle—from the monitor reading a new file to the AI making a decision and the actuator killing the process—took roughly 40 to 45 milliseconds. This speed is critical. It proves that a ransomware script barely has time to encrypt a single text document before the autonomous engine recognizes the math anomaly and terminates the execution. 

Furthermore, having the live web dashboard update instantly allows an IT team to monitor a whole building of computers simultaneously without having to do the hard work themselves; the AI cleans up the mess and simply reports the outcome to the dashboard.

\section{Weaknesses, Limitations, and Future Work}

Like any software project, we found a few areas that need improvement. By stripping 79 network variables down to just 6, we definitely lost some tiny details that could have bumped our accuracy to 95\%. Second, the way we kill background processes relies heavily on Windows operating system APIs. If we wanted to deploy this on Mac or Linux servers, we would have to completely rewrite the mitigation scripts to use POSIX commands instead. 

We also realize that modern malware authors are smart; they are developing ``slow and low'' attacks mathematically designed to stay just beneath standard monitoring thresholds. Our Random forest is generally resilient to this, but extremely advanced, slow-moving payload deployments operating across a scale of several months could evade detection simply because our engine analyzes immediate time-slices rather than long-term historical context.

For future development, we want to explore federated learning. If we deployed this engine across an entire corporate network, we could have individual laptops anonymously share their localized threat data with a master server. This would allow the Random Forest model to continuously update its own weights without needing manual database retraining.

\section{Conclusion}

Through this project, we proved that AI doesn't just have to sit in the cloud analyzing logs after a company has already been hacked. By combining system-level telemetry with a streamlined Random Forest classifier, we created an autonomous bouncer that lives right inside the computer it's trying to protect. 

We successfully achieved a realistic 92.41\% threat detection accuracy running directly against real-world intrusion data. More importantly, we showed how that intelligence can be directly mapped to actionable OS-level commands that kill threats instantly without human permission. By writing our own data synchronization buffers, analyzing and tuning false positive thresholds, and developing a live React telemetry dashboard, we took a theoretical machine learning concept and turned it into a fully functional endpoint security tool. This project moves the conversation past traditional signature-based antiviruses and shows that the future of personal defense relies on fast, behavioral algorithms that fight back autonomously.

\bibliographystyle{IEEEtran}
\begin{thebibliography}{20}

\bibitem{sarraf2025ai}
A. Sarraf and S. Pal, ``AI-Driven Autonomous Cloud Security for Threat Detection and Response,'' \emph{International Journal of Engineering Research \& Technology (IJERET)}, 2025.

\bibitem{yan2025ml}
X. Yan \emph{et al.}, ``Machine Learning and Deep Learning-Based Ransomware Detection: A Comprehensive Review,'' \emph{Elsevier -- Computer Security Review (CSR)}, 2025.

\bibitem{venckauskas2025zero}
A. Ven\v{c}kauskas \emph{et al.}, ``Zero-Day Ransomware Detection Using Static Feature-Based Machine Learning,'' \emph{MDPI Applied Sciences}, 2025.

\bibitem{lee2025ml}
J. Lee \emph{et al.}, ``Machine Learning-Based Detection of Ransomware in Encrypted Files Using FPE Patterns,'' \emph{MDPI Sensors}, 2025.

\bibitem{alhuwayshil2025risk}
A. Alhuwayshil \emph{et al.}, ``Risk-Aware Machine Learning for Ransomware Threat Prioritization,'' \emph{MDPI Sensors}, 2025.

\bibitem{fang2024honeyfile}
L. Fang \emph{et al.}, ``Honeyfile Anomaly Detection for Early Ransomware Mitigation,'' \emph{MDPI Mathematics}, 2024.

\bibitem{cen2024zeroshot}
Y. Cen \emph{et al.}, ``Zero-Shot Learning for Zero-Day Ransomware Early Detection,'' \emph{Elsevier Computers \& Security}, 2024.

\bibitem{aljabri2024memory}
M. Aljabri \emph{et al.}, ``Ransomware Detection Using Memory Dumps and Machine Learning,'' \emph{Elsevier JISA}, 2024.

\bibitem{kritika2024survey}
Kritika and R. Saxena, ``Review of AI-Based Ransomware Detection,'' \emph{Elsevier}, 2024.

\bibitem{jawad2024ml}
H. Jawad and H. Salman, ``Comparative Study of ML Models for Ransomware Analysis,'' \emph{IJ Safety \& Security Engg}, 2024.

\bibitem{toure2023hybrid}
A. Touré \emph{et al.}, ``Zero-Day Detection in Network Flows via Hybrid ML,'' \emph{Elsevier Info Sciences}, 2023.

\bibitem{shaukat2023systematic}
K. Shaukat \emph{et al.}, ``ML-based Ransomware Detection: Systematic Review,'' \emph{Springer}, 2023.

\bibitem{ring2023flow}
M. Ring \emph{et al.}, ``Network-Level Ransomware Detection via Flow-Based ML,'' \emph{IEEE CNS}, 2023.

\bibitem{homayoun2022behavioral}
S. Homayoun \emph{et al.}, ``Detecting Ransomware via Behavioral ML Analysis,'' \emph{IEEE BigData}, 2022.

\bibitem{conti2022security}
M. Conti \emph{et al.}, ``Malware \& Ransomware Trends Security Survey,'' \emph{Elsevier}, 2022.

\bibitem{lashkari2022deep}
A. Lashkari \emph{et al.}, ``Intelligent Intrusion Detection via Deep Learning IDS,'' \emph{Springer}, 2022.

\bibitem{vinayakumar2021deep}
R. Vinayakumar \emph{et al.}, ``Intrusion Detection Using Deep Neural Networks,'' \emph{Springer}, 2021.

\bibitem{scaife2021cryptolock}
N. Scaife \emph{et al.}, ``CryptoLock Defense for Practical Ransomware Mitigation,'' \emph{IEEE ICDCS}, 2021.

\bibitem{sood2021apt}
A. Sood and R. Enbody, ``Advanced Persistent Threat Analysis Fundamentals,'' \emph{IEEE S\&P}, 2021.

\bibitem{goodfellow2016deep}
I. Goodfellow, Y. Bengio, and A. Courville, ``Deep Learning,'' \emph{MIT Press}, 2016.

\end{thebibliography}

\end{document}
